\documentclass{aa}

\usepackage{graphicx}
\usepackage{amsmath}
\usepackage{booktabs}
\usepackage{float}
\usepackage{txfonts}

\newcommand{\compacttables}{\footnotesize\setlength{\tabcolsep}{3pt}}

%-------------------------------------------------
% DADOS DO RELATÓRIO
%-------------------------------------------------

\newcommand{\aluno}{Nome do Aluno}
\newcommand{\experimento}{Espectrofotometria estelar}

\newcommand{\disciplina}{Astronomia}
\newcommand{\professor}{Prof. Dr. Rafael C. R. de Lima}
\newcommand{\datarelatorio}{\today}

\begin{document}

\title{\experimento}
\subtitle{Relatório de laboratório em formato Astronomy \& Astrophysics}

\author{\aluno}

\institute{
Universidade do Estado de Santa Catarina, Centro de Ciências Tecnológicas -- CCT\\
\email{Disciplina: \disciplina; professor: \professor}
}

\date{\datarelatorio}

\abstract
{Este relatório apresenta uma análise espectrofotométrica de estrelas observadas no laboratório virtual.}
{O objetivo é extrair espectros a partir de aberturas fotométricas, subtrair o fundo local, ajustar um contínuo aproximado de corpo negro e medir grandezas fotométricas e linhas de absorção.}
{Foram registradas uma ou mais estrelas no simulador. Para cada alvo, o laboratório exportou o fluxo da abertura da estrela, o fundo escalado, o fluxo líquido, a fotometria sintética e as linhas de absorção marcadas.}
{Os resultados principais são apresentados nas tabelas e figuras geradas automaticamente. O aluno deve completar a interpretação física, comparando temperaturas, cores, fluxo integrado e possíveis identificações das linhas.}
{O formato segue a classe oficial \texttt{aa.cls} da revista Astronomy \& Astrophysics.}

\keywords{Stars: fundamental parameters -- Techniques: spectroscopic -- Techniques: photometric -- Methods: data analysis}

\maketitle

%=================================================
\section{Introdução}

Descreva aqui o objetivo do experimento e explique por que espectros estelares permitem inferir temperatura, cor, fluxo relativo e composição atmosférica.

%=================================================
\section{Fundamentação teórica}

O contínuo espectral de uma estrela pode ser aproximado, em primeira ordem, por uma curva de corpo negro. A posição do máximo da emissão está relacionada à temperatura pela lei de Wien,

\begin{equation}
\lambda_{\max}T = 2{,}898\times 10^{-3}\ {\rm m\,K}.
\end{equation}

A fotometria sintética em bandas permite construir índices de cor. Neste laboratório, usamos

\begin{equation}
B - V = -2{,}5\log_{10}\left(\frac{F_B}{F_V}\right).
\end{equation}

Linhas de absorção deslocadas em comprimento de onda podem ser interpretadas, na aproximação não relativística, por

\begin{equation}
\frac{v_r}{c}\simeq \frac{\lambda_{\rm obs}-\lambda_0}{\lambda_0}.
\end{equation}

%=================================================
\section{Observações e redução dos dados}

As tabelas e figuras abaixo foram exportadas diretamente do laboratório após o registro das estrelas observadas. Para cada estrela, o fluxo líquido corresponde ao fluxo medido na abertura da estrela menos o fundo local escalado pela área da abertura.

\input{tabelas/tabela_resumo_observacoes.tex}

\input{tabelas/tabela_fotometria.tex}

\input{tabelas/tabela_linhas_absorcao.tex}

\input{secoes/observacoes.tex}

%=================================================
\section{Análise}

Apresente aqui os cálculos feitos com os dados exportados. Use os arquivos em \texttt{dados/} para consultar os valores de fluxo por comprimento de onda de cada estrela.

Compare as temperaturas obtidas pelo ajuste de corpo negro, os valores de $B-V$, o fluxo integrado e as linhas de absorção marcadas. Discuta quais estrelas são mais quentes, quais são mais frias e se alguma estrela é semelhante ao Sol.

%=================================================
\section{Discussão}

Discuta as limitações do método: escolha da abertura da estrela, escolha da abertura de fundo, ruído estatístico, linhas de absorção sobrepostas ao contínuo e intervalo de comprimento de onda usado no ajuste.

%=================================================
\section{Conclusões}

Escreva aqui as principais conclusões do experimento.

\end{document}
